\section*{Streszczenie}

Celem pracy była analiza porównawcza architektury monolitycznej oraz
architektury mikroserwisowej pod kątem ich wydajności, skalowalności
oraz kosztów utrzymania w środowiskach chmurowych. W dobie rosnącej
złożoności systemów informatycznych oraz dynamicznego przyrostu liczby
użytkowników coraz większe znaczenie zyskują modele architektoniczne
umożliwiające efektywne zarządzanie zasobami, elastyczność oraz
niezawodność. W pracy zbadano, w jaki sposób wybór architektury wpływa
na realne parametry działania aplikacji, a także jakie konsekwencje
niesie dla organizacji pod względem operacyjnym i finansowym.

Przeprowadzony przegląd literatury pozwolił zidentyfikować aktualne
podejścia do pomiaru wydajności i skalowalności systemów rozproszonych
oraz istniejące narzędzia benchmarkingowe. Jednak większość dostępnych
badań analizuje architektury w oderwaniu od praktycznego kontekstu
wdrożenia w chmurze, co wskazało na istnienie luki badawczej — brak
kompleksowego porównania monolitu i mikroserwisów uruchomionych w
porównywalnym środowisku, z zastosowaniem jednolitych metod obciążenia
i identycznych scenariuszy testowych.

W ramach pracy zaimplementowano aplikację benchmarkingową w dwóch
wariantach architektonicznych — monolitycznym oraz mikroserwisowym —
oraz wdrożono oba modele na platformie chmurowej. Zastosowano testy
obciążeniowe generowane za pomocą dedykowanego narzędzia oraz
przeprowadzono analizę skalowania pionowego i poziomego. Dodatkowo
zbadano wpływ konfiguracji infrastruktury, mechanizmów komunikacji,
wielowątkowości oraz sposobu zarządzania zasobami na ogólną
efektywność aplikacji.

Wyniki badań wskazują, że architektura monolityczna charakteryzuje się
wyższą wydajnością przy niewielkim obciążeniu oraz niższymi kosztami
utrzymania w prostych instalacjach, natomiast architektura
mikroserwisowa oferuje znacząco lepszą skalowalność poziomą,
niezależność komponentów oraz większą elastyczność w zakresie
optymalizacji wybranych obszarów systemu. Jednocześnie mikroserwisy
wprowadzają koszty związane z komunikacją między usługami,
złożonością infrastruktury oraz większym nakładem na utrzymanie.

Praca kończy się omówieniem ograniczeń badania, wkładu pracy w obszar
inżynierii oprogramowania oraz propozycją kierunków dalszych badań,
obejmujących m.in.\ analizę architektur hybrydowych,
automatyczne skalowanie zasobów oraz wpływ technologii serverless na
koszty operacyjne aplikacji.

\vspace{0.5cm}

\noindent \textbf{Słowa kluczowe:} Informatyka, Mikroserwisy, Monolit, Analiza kosztów, GCP, gRPC, Aplikacja webowa
\vspace{0.5cm}

\noindent \textbf{Dziedzina nauki i techniki, zgodnie z wymogami OECD:} Nauki o komputerach i informatyka
\vspace{0.5cm}